% Based on templates/lecture-example.tex.
% Original worked examples: IO and Disk.pdf, slides 10.17--10.19.
\begingroup
\renewcommand{\recordid}[1]{}
\renewcommand{\relatedtopic}[2]{\hyperref[topic:#1]{#2}}
\exercisegroup{SC2005-L22-EX}{第 22 讲：磁盘随机与顺序读取时间}

\examplemeta
  {SC2005-L22-EX}
  {2026-11-05}
  {IO and Disk 10.17--10.19（PDF pp.~17--19）}
  {按讲义原例题整理，补全换算及定位条件}
  {已完成学习并要求整理}

\linespread{1.15}\selectfont
\setlength{\parskip}{0.4em}

\exercisequestion{SC2005-L22-E01}{分散扇区的随机读取}
\textbf{题意与来源。}讲义 10.17--10.18：平均寻道时间 10 ms，转一圈 8 ms，平均旋转等待 4 ms；每扇区 512 B，每磁道 1024 个扇区。先求随机读一个扇区的平均时间，再求读取总共 4 KB、128 KB、1 MB 的时间。\textbf{每个扇区分别随机定位}，各次访问依次完成；沿用课件的二进制单位，$1\,\mathrm{KB}=1024$ B、$1\,\mathrm{MB}=1024^2$ B。对应\relatedtopic{SC2005-L22-T04}{磁盘读写时间}。

\begin{workedsolution}
单个扇区需要一次寻道、一次旋转等待及一个扇区的传输：
\[
  T_1=10+4+\frac{1}{1024}\times8
  =14.0078125\text{ ms}\approx14\text{ ms}.
\]
总扇区数为数据字节数除以 512，独立访问 $N$ 次需 $NT_1$。按讲义取每次约 14 ms：
\begin{center}
\begin{tabular}{@{}lll@{}}
\toprule
数据量 & 扇区数 $N$ & 总时间（按 $T_1\approx14$ ms） \\
\midrule
4 KB & $4\times1024/512=8$ & $8\times14=112$ ms \\
128 KB & $128\times1024/512=256$ & $256\times14=3584$ ms $=3.584$ s \\
1 MB & $1024^2/512=2048$ & $2048\times14=28672$ ms $=28.672$ s \\
\bottomrule
\end{tabular}
\end{center}
保留单扇区传输时间时，三项分别为 112.0625 ms、3.586 s、28.688 s。主要成本来自重复定位；连续排列的多个扇区无需每个都单独定位，应按其布局计算。
\end{workedsolution}

\exercisequestion{SC2005-L22-E02}{连续数据的顺序读取}
\textbf{题意与来源。}讲义 10.17、10.19：使用相同参数，即寻道 10 ms、整圈旋转 8 ms、平均旋转等待 4 ms、扇区 512 B、每磁道 1024 个扇区，估算连续读取 4 KB、128 KB、1 MB 的时间。KB、MB 仍按二进制换算；前两项假定数据全部落在同一磁道；1 MB 按占满两条完整磁道、每磁道各付一次平均定位时间计算。对应\relatedtopic{SC2005-L22-T04}{磁盘读写时间}。

\begin{workedsolution}
\textbf{4 KB 与 128 KB。}一次平均定位为 $10+4=14$ ms，只在其后增加连续数据的传输时间：
\[
  T_{4\mathrm{KB}}=14+\frac{8}{1024}\times8
  =\boxed{14.0625\text{ ms}},\qquad
  T_{128\mathrm{KB}}=14+\frac{256}{1024}\times8
  =\boxed{16\text{ ms}}.
\]
数据量增大而定位次数不变，因此增加的主要是传输时间。

\textbf{1 MB。}每磁道容量为 $1024\times512=524288$ B，即 512 KiB。1 MiB 占两条完整磁道，每条磁道需 $14+8=22$ ms，故
\[
  T_{1\mathrm{MB}}=2(14+8)=\boxed{44\text{ ms}}.
\]
44 ms 采用讲义的两磁道定位模型。若起点未对齐磁道边界，1 MiB 可跨三条磁道；实际磁道切换成本需按硬件或题目条件给定。

\textbf{比较。}分散读取的每个扇区都付定位成本，顺序读取则让一次定位服务更多数据。计算前先确定扇区数、跨越磁道数和定位次数，不能仅凭请求大小决定耗时。
\end{workedsolution}
\endgroup
